\section{Chat Proving Ground 与 API Contract}

最后一节把研究与产品连接起来。Chat experience 由供应商控制 UI、system prompt、tool、rollback 与 user education，可以快速试验；API 一旦被客户写进 production，schema、model behavior、latency、error、pricing 和 retirement 都成为外部 contract。正如课堂所说，“APIs are forever”不是字面永不变化，而是变化成本落在大量客户系统上。

\subsection{从实验到 versioned primitive}

Chat 中验证 PDF upload、long context 或 tool behavior 后，团队需要把隐式 UX 约束转成 API schema、auth、quota、idempotency、error model、version 和 docs。\term{API deprecation}（弃用）是宣布某模型或参数不再推荐，并提供 replacement、migration guidance 与 retirement date；retirement 后请求失败，因此客户必须在窗口内测试和迁移。

\lecturefigure{16-chat-to-api.png}{Chat 负责快速试验，API 则把成熟能力发布为长期兼容 contract。}{本地字幕 39:00--41:56；Anthropic API lifecycle docs；概念重绘。}

\begin{knowledgebox}{读图：Release 之前要补齐哪些证据}
Chat experiment 先观察 usefulness、安全和 failure；evidence 形成 acceptance criteria；API release 固化 schema/version/SLO；lifecycle 提供 active、legacy、deprecated、retired 状态与 migration。客户 continuity 可能比“最新模型更聪明”更重要，因此 replacement 必须在真实 workload 上回归，而不能只看 benchmark。
\end{knowledgebox}

\teachervoice{Mann 用 Claude 旧模型仍在某处生产环境说明 API inertia：客户可能没有工程资源迁移，也可能把旧行为写进合规流程。课堂提醒 frontier lab，发布 API 就等于接受 support、compatibility 和 deprecation 责任。}

\begin{warningbox}{模型升级不是普通 binary upgrade}
新模型可能改变 refusal、format、tool call、latency、token usage 与 nondeterministic distribution。Migration 需要 golden tasks、safety regression、cost/SLO test、shadow traffic、rollback 和 customer communication；同名 alias 静默切换会让事故难以归因。
\end{warningbox}

\subsection{本章小结}

Chat 与 API 的速度不同是合理的。Proving ground 帮助发现行为，API contract 要求版本化、可测试和可迁移；产品纪律是 frontier capability 成为基础设施的最后一步。

